% Technical note: EqTide Darwin-Kaula constituent decomposition
% Compile: pdflatex tidal_constituents.tex

\documentclass[11pt]{article}

\usepackage[margin=1in]{geometry}
\usepackage{amsmath}
\usepackage{amssymb}
\usepackage{booktabs}
\usepackage{graphicx}
\usepackage[colorlinks=true,linkcolor=blue,urlcolor=blue]{hyperref}

\newcommand{\eqtide}{\textsc{EqTide}}
\newcommand{\Om}{\Omega}

\title{Resolving the Darwin--Kaula Tidal Constituents in \eqtide:\\
Frequencies, Amplitudes, and Dissipated Power}
\author{Rory Barnes\\\small University of Washington}
\date{7 August 2026}

\begin{document}
\maketitle

\begin{abstract}
\noindent
\eqtide's constant-phase-lag (CPL) tidal model expands the degree-2
tide-raising potential into six discrete Darwin--Kaula harmonics, but it has
only ever used the \emph{sign} of each harmonic's frequency, since that is all
the secular equations of motion require. Everything else about the individual
constituents was computed internally and discarded. This note documents a
modification that recovers it and reports, for each of the six constituents,
its forcing frequency, its tide height, and the power it dissipates. These are
the quantities physical oceanography works with, and exposing them is a
prerequisite for coupling an orbital evolution model to an ocean-response
model --- the motivating application being Enceladus. Alongside the
decomposition we introduce the displacement Love number $h_2$, which \eqtide\
has never carried; without it the code can report the tide of the raising
potential but not the tide a gauge would measure. The decomposition of the
power is exact, being a regrouping of terms already present, and neither it
nor $h_2$ alters any orbital or rotational evolution. We validate in two
opposite dynamical limits. For the fast-rotating Earth the six constituents
land on named solar Doodson lines (S$_2$, T$_2$, R$_2$, P$_1$, K$_1$,
S$_\mathrm{a}$), with periods matching the tide tables to $\sim\!10^{-7}$
relative and within-species amplitude ratios matching the Cartwright--Tayler
expansion to a few parts in $10^3$. For synchronously locked Enceladus the
constituent structure collapses onto a single frequency, the heating
partitions in exact rational proportions $7/8$, $3/28$, $1/56$, and the summed
constituent power reproduces the textbook closed form
$\tfrac{21}{2}(k_2/Q)GM_p^2R^5ne^2/a^6$ to $1.4\times10^{-7}$.
\textbf{This is preliminary work circulated for independent verification};
\S\ref{sec:caveats} lists the points where a second opinion would be most
valuable.
\end{abstract}

\section{Motivation and scope}

Physical oceanography describes a tide as a superposition of harmonic
constituents, each with its own frequency, amplitude, and phase. Tidal
prediction, ocean-resonance analysis, and dissipation budgets in subsurface
oceans are all organized around that decomposition. \eqtide, by contrast, is
built to evolve orbits and spins: it needs only the secular, orbit-averaged
torques and the total dissipated power, so its representation of the tide has
never been expressed in constituent form.

The two descriptions are not in conflict --- \eqtide's CPL model \emph{is} a
constituent model internally. The purpose of this modification is to make that
structure explicit, so that a run can hand an ocean model a list of
(frequency, amplitude) pairs rather than a single bulk heating rate.

Scope is deliberately narrow. The new quantities are diagnostics: they are
derived from the model state at the moment they are reported, and nothing
computed here feeds back into the equations of motion. No orbit, spin, or
obliquity evolves differently than it did before. The one genuinely new
physical ingredient is the Love number $h_2$ (\S\ref{sec:h2}), and it too
enters only the reported tide heights.

\section{The CPL model and its constituents}

\subsection{Kaula expansion of the tide-raising potential}

The potential raised on a body of radius $R$ by a perturber of mass $M_p$ at
semi-major axis $a$ is expanded, at degree $l=2$, in the Kaula form
\begin{equation}
  U = \frac{GM_p}{a}\sum_{m=0}^{2}\sum_{p=0}^{2}\sum_{q=-\infty}^{\infty}
  \left(\frac{R}{a}\right)^{2} F_{2mp}(\psi)\,G_{2pq}(e)\,
  \cos\!\big(\sigma_{mpq}t + \phi_{mpq}\big),
\end{equation}
where $\psi$ is the obliquity, $e$ the eccentricity, $F_{2mp}$ the inclination
functions and $G_{2pq}$ the eccentricity functions. Each harmonic is forced at
\begin{equation}
  \sigma_{mpq} = (2-2p+q)\,n - m\,\Om ,
  \label{eq:sigma}
\end{equation}
with $n$ the mean motion and $\Om$ the rotation rate.

In a constant-phase-lag model each harmonic is assigned a lag
$\varepsilon_{mpq}$ whose \emph{magnitude} is fixed by the quality factor $Q$
and whose \emph{sign} follows the sign of $\sigma_{mpq}$:
\begin{equation}
  \varepsilon_i = \operatorname{sgn}(\sigma_i),
  \qquad
  \varepsilon_i \in \{-1,0,+1\}.
  \label{eq:eps}
\end{equation}
The zero case is not a rounding artifact but a physical statement: a harmonic
at exactly zero frequency is a permanent, unlagged bulge. It matters at
Enceladus (\S\ref{sec:enceladus}).

\subsection{The six retained harmonics}

Truncated at $\mathcal{O}(e)$ in amplitude and first order in the inclination
functions --- the order CPL itself retains --- six harmonics survive.
Table~\ref{tab:constituents} lists them with the names used throughout this
note. Doodson analogues are for the Earth--Sun system only and are \emph{not}
general (see \S\ref{sec:enceladus}).

\begin{table}[h]
\centering
\caption{The six Darwin--Kaula constituents resolved by \eqtide's CPL model.}
\label{tab:constituents}
\begin{tabular}{llccl}
\toprule
Name & $(m,p,q)$ & Frequency $\sigma$ & $G_{2pq}$ & Earth--Sun \\
\midrule
Semi-diurnal        & $(2,0,0)$  & $2\Om-2n$ & $1-\tfrac52 e^2$ & S$_2$ \\
Ecc.\ upper         & $(2,0,+1)$ & $2\Om-3n$ & $\tfrac72 e$     & T$_2$ \\
Ecc.\ lower         & $(2,0,-1)$ & $2\Om-n$  & $-\tfrac12 e$    & R$_2$ \\
Radial (zonal)      & $(0,1,+1)$ & $n$       & $\tfrac32 e$     & S$_\mathrm{a}$* \\
Obliquity, diurnal  & $(1,0,0)$  & $\Om-2n$  & $\simeq 1$       & P$_1$ \\
Obliquity, sidereal & $(1,1,0)$  & $\Om$     & $\simeq 1$       & K$_1$ \\
\bottomrule
\end{tabular}
\end{table}

Frequencies are reported as magnitudes, $|\sigma_i|$: a tidal frequency is
positive by convention, and the sign of the forcing argument is already
carried by $\varepsilon_i$.

\section{What is computed}

\subsection{Frequency}

Directly from Eq.~\eqref{eq:sigma}, evaluated at the body's instantaneous
$\Om$ and the orbit's $n$. This information already existed inside the model
and was being discarded by the $\operatorname{sgn}$ in Eq.~\eqref{eq:eps};
recovering it required no new physics.

\subsection{The tide of the raising potential}

The degree-2 potential of constituent $i$ has amplitude
$\Psi_i = (GM_pR^2/a^3)\,A_i$, so the equilibrium surface displacement,
$\zeta_i = \Psi_i/g$ with $g = GM_b/R^2$, is
\begin{equation}
  \boxed{\;
  \zeta_i = \frac{M_p}{M_b}\,\frac{R^4}{a^3}\,A_i \;}
  \qquad
  A_i \equiv \tfrac{1}{3}F_{2mp}(\psi)\,G_{2pq}(e).
  \label{eq:amp}
\end{equation}
The factor $1/3$ normalizes every inclination function so that the principal
semi-diurnal term reduces to unity at $\psi=0$, $e=0$; this makes
Eq.~\eqref{eq:amp} return the textbook equilibrium tide height directly and
leaves all constituent ratios, within and across species, equal to those of
the unnormalized Kaula expansion. The normalized inclination functions are
\begin{align}
  \tfrac13 F_{220} &= \cos^4(\psi/2), &
  \tfrac13 F_{201} &= \tfrac12\sin^2\psi - \tfrac13, \\
  \tfrac13 F_{210} &= \tfrac12\sin\psi\,(1+\cos\psi), &
  \tfrac13 F_{211} &= \sin\psi\cos\psi .
\end{align}

Equation~\eqref{eq:amp} is the tide of the \emph{raising potential}. It is not
yet the tide anything would measure, which is the subject of the next section.

\subsection{Which tide? The role of \texorpdfstring{$h_2$}{h2}}
\label{sec:h2}

Two Love numbers of degree 2 are needed to turn a forcing potential into an
observable tide, and \eqtide\ has historically carried only one of them.

\begin{itemize}
\item $k_2$, the \emph{potential} Love number, is the ratio of the additional
gravitational potential generated by the deformed body to the raising
potential. It is what sets the tidal torque and, through $\operatorname{Im}
(k_2) = -k_2/Q$, the dissipation.
\item $h_2$, the \emph{displacement} Love number, is the ratio of the radial
displacement of the solid surface to $\zeta$. It sets how far the ground
itself moves.
\end{itemize}

The distinction matters because an ocean is measured against a sea floor that
is moving too. Three heights follow from the same forcing $\zeta$:
\begin{align}
  \text{solid surface rises by} \quad & h_2\,\zeta, \\
  \text{geoid rises by} \quad & (1+k_2)\,\zeta, \\
  \text{water depth changes by} \quad & \gamma_2\,\zeta,
  \qquad \gamma_2 \equiv 1 + k_2 - h_2 .
  \label{eq:gamma}
\end{align}
The middle line includes the self-attraction of the redistributed mass, which
lifts the equipotential a further $k_2\zeta$ above the raising value. An ocean
fills to the geoid; a tide gauge measures the gap between the geoid and the
solid surface beneath it. The difference, $\gamma_2$, is the classical
diminishing factor. For Earth, with $k_2 = 0.298$ and $h_2 = 0.6103$,
\begin{equation}
  \gamma_2 = 1 + 0.298 - 0.6103 = 0.6877,
\end{equation}
so the ocean tide is about $31\%$ smaller than the naive equilibrium tide of
the raising potential. Reporting $\zeta$ alone would overstate every tide
height by that factor, which is why $h_2$ had to be added before these
diagnostics could be handed to an oceanographer.

\paragraph{What $h_2$ does \emph{not} affect.} It is worth being explicit,
because it bounds the risk of the change. The dissipation is governed by
$\operatorname{Im}(k_2)$; $h_2$ appears nowhere in the energy budget, nowhere
in Eq.~\eqref{eq:sigma}, and nowhere in the equations of motion. Since
$\gamma_2$ is also the same for every constituent, applying it rescales the
whole set uniformly and changes no amplitude \emph{ratio}. Introducing $h_2$
therefore changes reported tide heights and nothing else --- a claim we verify
numerically in \S\ref{sec:earth}.

\paragraph{Default value.} When $h_2$ is not supplied we use
\begin{equation}
  h_2 = \tfrac53 k_2,
\end{equation}
which is exact for a homogeneous incompressible elastic sphere: with effective
rigidity $\tilde\mu$, $k_2 = (3/2)/(1+\tilde\mu)$ and
$h_2 = (5/2)/(1+\tilde\mu)$. That is the same idealization a single global
$k_2$ already assumes, which makes it the self-consistent fallback. It should
not be mistaken for a measurement: for Earth it gives $h_2 = 0.50$ against the
observed $0.61$, an $18\%$ error in $h_2$ and a $16\%$ error in $\gamma_2$.
Where $h_2$ is known it should be set explicitly.

\subsection{Dissipated power, and why the split is exact}

\eqtide\ computes the total tidal power as a sum of an orbital and a
rotational term,
\begin{equation}
  P = \frac{Z}{8}\left[\frac{\Om}{n}\,\Gamma_{\rm rot} - \Gamma_{\rm orb}\right],
  \label{eq:ptot}
\end{equation}
where $\Gamma_{\rm orb}$ and $\Gamma_{\rm rot}$ are each a linear combination
of the $\varepsilon_i$ and
\begin{equation}
  Z = 3\,\frac{k_2}{Q}\,
      \frac{G^2 M_p^2 (M_b+M_p) R^5}{a^9\,n}.
  \label{eq:Z}
\end{equation}
Written out, the two combinations are
\begin{align}
  \Gamma_{\rm orb} &= 4\varepsilon_0
    + e^2\!\left(-20\varepsilon_0 + \tfrac{147}{2}\varepsilon_1
      + \tfrac12\varepsilon_2 - 3\varepsilon_5\right)
    - 4\sin^2\!\psi\,(\varepsilon_0 - \varepsilon_8), \\
  \Gamma_{\rm rot} &= 4\varepsilon_0
    + e^2\!\left(-20\varepsilon_0 + 49\varepsilon_1 + \varepsilon_2\right)
    + 2\sin^2\!\psi\,(-2\varepsilon_0 + \varepsilon_8 + \varepsilon_9),
\end{align}
where the subscript on $\varepsilon$ labels the constituent in the order of
Table~\ref{tab:constituents} ($0$ semi-diurnal, $1$ and $2$ the eccentricity
sidebands, $5$ radial, $8$ and $9$ the obliquity terms). Both are already sums
over constituents. Collecting the coefficient of each $\varepsilon_i$
separately, rather than summing them, gives the power carried by constituent
$i$:
\begin{equation}
  \boxed{\;
  P_i = \frac{Z}{8}\,\varepsilon_i
        \left[\frac{\Om}{n}\,w^{\rm rot}_i - w^{\rm orb}_i\right] \;}
  \label{eq:pi}
\end{equation}
with the weights of Table~\ref{tab:weights}. Because Eq.~\eqref{eq:pi} is a
regrouping of the same terms, $\sum_i P_i \equiv P$ identically --- not
approximately. This is a strong internal check, and we exercise it in both
validation cases below, where $P$ is obtained by a completely independent
route.

\begin{table}[h]
\centering
\caption{Weights entering Eq.~\eqref{eq:pi}, read term-by-term out of
$\Gamma_{\rm orb}$ and $\Gamma_{\rm rot}$.}
\label{tab:weights}
\begin{tabular}{lcc}
\toprule
Constituent & $w^{\rm orb}_i$ & $w^{\rm rot}_i$ \\
\midrule
Semi-diurnal        & $4 - 20e^2 - 4\sin^2\psi$ & $4 - 20e^2 - 4\sin^2\psi$ \\
Ecc.\ upper         & $\tfrac{147}{2}e^2$       & $49e^2$ \\
Ecc.\ lower         & $\tfrac12 e^2$            & $e^2$ \\
Radial              & $-3e^2$                   & $0$ \\
Obliquity, diurnal  & $4\sin^2\psi$             & $2\sin^2\psi$ \\
Obliquity, sidereal & $0$                       & $2\sin^2\psi$ \\
\bottomrule
\end{tabular}
\end{table}

\subsection{Applicability}

The constituent decomposition is defined only for the constant-phase-lag
model. A constant-time-lag rheology responds continuously across frequency and
so has no discrete set of constituents to report; asking for one there returns
a not-applicable sentinel of $-1$, which is unambiguous because a frequency,
an amplitude and a dissipated power are all non-negative by construction. The
Love numbers and $\gamma_2$, being properties of the body's elastic response
rather than of the tidal model, remain meaningful and are reported in every
case. Note that $\gamma_2$ could not use the same sentinel: it legitimately
reaches zero for a body in the fluid limit, where $h_2 = 1+k_2$ exactly and a
freely flowing body raises no tide relative to its own deforming surface.

\section{Validation I: Earth--Sun, the fast-rotator limit}
\label{sec:earth}

Earth rotates $\sim\!365$ times per orbit, so $\Om \gg n$, every
$\varepsilon_i = +1$, and the six constituents separate cleanly into three
species. This is the regime the classical tide tables were built for, so it
admits a direct external check.

\paragraph{Configuration.} Earth: $M=M_\oplus$, $R=R_\oplus$,
$P_{\rm rot}=0.99726968$~d (sidereal), $\psi=23.439281^\circ$, $e=0.0167086$,
$a=1.00000102$~au, $k_2=0.298$, $h_2=0.6103$, $Q=12$ (ocean-dominated). The
Sun is the sole perturber. \textbf{There is no Moon in this configuration},
which is why no M$_2$ appears: every line below is solar.

\begin{table}[h]
\centering
\caption{Earth--Sun constituent periods and tide heights. ``Reference'' is the
tabulated Doodson period; ``raising'' is $\zeta$ from Eq.~\eqref{eq:amp} and
``ocean'' is $\gamma_2\zeta$ from Eq.~\eqref{eq:gamma}.}
\label{tab:earth}
\small
\begin{tabular}{llrrrrr}
\toprule
Constituent & Doodson & Period (h) & Reference (h) & Rel.\ err &
Raising (m) & Ocean (m) \\
\midrule
Semi-diurnal & S$_2$ & 12.000000 & 12.000000 & $4.3\times10^{-9}$ & 0.151168 & 0.103959 \\
Ecc.\ upper  & T$_2$ & 12.016449 & 12.016449 & $3.0\times10^{-8}$ & 0.008847 & 0.006084 \\
Ecc.\ lower  & R$_2$ & 11.983596 & 11.983596 & $-2.4\times10^{-8}$ & 0.001264 & 0.000869 \\
Obl., diurnal & P$_1$ & 24.065888 & 24.065890 & $-9.7\times10^{-8}$ & 0.062763 & 0.043162 \\
Obl., sidereal & K$_1$ & 23.934472 & 23.934470 & $9.7\times10^{-8}$ & 0.060062 & 0.041305 \\
Radial & S$_\mathrm{a}$* & 8766.165025 & 8766.152700 & $1.4\times10^{-6}$ & 0.001049 & 0.000721 \\
\bottomrule
\end{tabular}
\end{table}

Five of the six identifications are exact. \textbf{S$_\mathrm{a}$* is a loose
match and is flagged as such}: \eqtide\ places this line at the Keplerian mean
motion, i.e.\ the \emph{sidereal} year (365.256363~d), whereas Doodson's
S$_\mathrm{a}$ is defined on the \emph{tropical} year (365.2422~d). The two
differ by the precession of the equinoxes, which \eqtide\ does not model. The
reference quoted is therefore the sidereal year, and the residual
$1.4\times10^{-6}$ is not a physical discrepancy.

The semi-diurnal ocean tide of 0.104~m is the right size for the solar
equilibrium tide, and is the number an ocean model should be given; the
0.151~m raising value is not.

\paragraph{Amplitude ratios.} Amplitudes are compared as ratios \emph{within}
each species, because the harmonic tables carry a different geodetic
normalization for the long-period, diurnal, and semi-diurnal species; a
cross-species ratio would be meaningless. Table~\ref{tab:earthratios} compares
against Cartwright \& Tayler (1971). Agreement is a few parts in $10^3$,
consistent with the truncation of the eccentricity functions at
$\mathcal{O}(e)$ plus rounding in the published tables. These ratios are
identical whether formed from the raising or the ocean column, which is the
numerical statement that $\gamma_2$ is constituent-independent.

\begin{table}[h]
\centering
\caption{Earth--Sun within-species amplitude ratios against the
Cartwright \& Tayler (1971) expansion.}
\label{tab:earthratios}
\begin{tabular}{lrrr}
\toprule
Ratio & \eqtide & Cartwright--Tayler & Rel.\ err \\
\midrule
T$_2$/S$_2$ & 0.058521 & 0.058459 & $1.1\times10^{-3}$ \\
R$_2$/S$_2$ & 0.008360 & 0.008372 & $-1.4\times10^{-3}$ \\
K$_1$/P$_1$ & 0.956965 & 0.958616 & $-1.7\times10^{-3}$ \\
\bottomrule
\end{tabular}
\end{table}

\paragraph{Power closure.} Table~\ref{tab:earthpower} gives the power budget.
The six constituent powers close against the independently computed total to
$2.4\times10^{-9}$ relative, confirming the exactness of Eq.~\eqref{eq:pi} in
the fast-rotator regime. The budget is numerically unchanged by the
introduction of $h_2$, as \S\ref{sec:h2} requires.

\begin{table}[h]
\centering
\caption{Earth--Sun power budget. The six constituent powers sum to the total
tidal power, which is obtained by an independent route.}
\label{tab:earthpower}
\begin{tabular}{lr}
\toprule
Constituent & Power (TW) \\
\midrule
S$_2$ & 0.56570142 \\
T$_2$ & 0.00229897 \\
R$_2$ & 0.00004705 \\
P$_1$ & 0.05310941 \\
K$_1$ & 0.05340102 \\
S$_\mathrm{a}$* & 0.00000039 \\
\midrule
Sum & 0.67455825 \\
Independent total & 0.67455825 \\
Residual & $2.4\times10^{-9}$ \\
\bottomrule
\end{tabular}
\end{table}

\section{Validation II: Enceladus--Saturn, the synchronous limit}
\label{sec:enceladus}

Enceladus is the opposite dynamical limit and the case of interest. It is
synchronously locked, $\Om = n$, and the constituent structure
\emph{collapses}.

\paragraph{Configuration.} Enceladus: $M=1.0802\times10^{20}$~kg,
$R=252.1$~km, $\psi=0.0005^\circ$, $e=0.0047$, $a=237{,}948$~km, $k_2=0.1$,
$Q=20$. Saturn: $M=5.6834\times10^{26}$~kg. Rotation is held at the
equilibrium (synchronous) rate. Here $h_2$ is left at its default,
$\tfrac53 k_2 = 0.1667$, giving $\gamma_2 = 0.9333$ --- a number we report but
deliberately do not use, for reasons given below.

\paragraph{Frequency collapse.} Substituting $\Om=n$ into
Table~\ref{tab:constituents}: the semi-diurnal frequency $2\Om-2n$ goes to
\emph{exactly} zero, and all five remaining constituents land on $n$.

\begin{table}[h]
\centering
\caption{Enceladus constituents. The orbital period is 32.8923~h. Heights are
those of the raising potential, Eq.~\eqref{eq:amp}.}
\label{tab:enceladus}
\begin{tabular}{llrrr}
\toprule
Constituent & Character & Period (h) & Height (m) & Power (GW) \\
\midrule
Semi-diurnal        & permanent & $\infty$ & 1577.3398 & 0.000000 \\
Ecc.\ upper         & ecc $+$   & 32.8923 & 25.9487 & 6.512370 \\
Ecc.\ lower         & ecc $-$   & 32.8923 & 3.7070  & 0.132910 \\
Radial              & radial    & 32.8923 & 3.7070  & 0.797430 \\
Obliquity, diurnal  & obliquity & 32.8923 & 0.0138  & 0.000002 \\
Obliquity, sidereal & obliquity & 32.8923 & 0.0138  & 0.000002 \\
\midrule
\multicolumn{4}{l}{Total} & 7.442710 \\
\multicolumn{4}{l}{Sum of constituents} & 7.442714 \\
\bottomrule
\end{tabular}
\end{table}

\paragraph{There is no M$_2$/S$_2$ doublet here, and there cannot be.} This
point is worth stating plainly for an oceanographic reader. The semi-diurnal
splitting familiar from Earth arises from a \emph{fast} rotator sampling
\emph{two} perturbers. Enceladus is a slow synchronous rotator with one
dominant perturber; everything time-variable happens at a single frequency.
Any ocean-response model built on an Earth-like constituent list will be
mis-specified for a synchronous satellite.

\paragraph{The permanent tide does no work.} At exactly zero frequency,
Eq.~\eqref{eq:eps} gives $\varepsilon_0=0$, so by Eq.~\eqref{eq:pi} the
semi-diurnal constituent carries \emph{no} dissipation. Its 1577~m height is
the static tidal bulge of the raising potential, not an oscillation --- the
right order of magnitude, given that Enceladus is measurably triaxial at the
kilometre level.

\paragraph{The heating splits in exact rationals.} With $\Om=n$ and
$\psi\to0$, the lags are $\varepsilon_0=0$, $\varepsilon_1=-1$,
$\varepsilon_2=+1$, $\varepsilon_5=+1$, $\varepsilon_8=-1$,
$\varepsilon_9=+1$. Equation~\eqref{eq:pi} then gives, in units of $Ze^2/8$,
\begin{equation}
  P_{\rm ecc+} = 24.5, \qquad
  P_{\rm radial} = 3, \qquad
  P_{\rm ecc-} = 0.5, \qquad
  \sum_i P_i = 28,
\end{equation}
so the shares are exactly
\begin{equation}
  \frac{24.5}{28} = \frac{7}{8}, \qquad
  \frac{3}{28}, \qquad
  \frac{0.5}{28} = \frac{1}{56},
\end{equation}
i.e.\ $87.5000\%$, $10.7142\%$, $1.7858\%$ --- reproduced numerically to the
precision of the reported values. The eccentricity tide does essentially all
the work; the obliquity tides are negligible for any plausible Enceladus
obliquity.

\paragraph{The total reproduces the textbook closed form.} Setting
$\sum_i P_i = \tfrac{28}{8}Ze^2 = \tfrac72 Ze^2$, substituting
Eq.~\eqref{eq:Z}, and eliminating $G(M_b+M_p) = n^2a^3$ via Kepler's third law
gives
\begin{equation}
  P = \frac{21}{2}\,\frac{k_2}{Q}\,
      \frac{G M_p^2 R^5 n e^2}{a^6},
  \label{eq:closed}
\end{equation}
the standard result for a synchronous satellite.
Table~\ref{tab:enceladusclosed} gives the numerical comparison. This is a
genuine external check, not a tautology: the factor $21/2$ emerges from the
CPL phase-lag weights summed over constituents, while the comparison value is
built from the physical parameters alone. The residual is limited by the
precision of the reported inputs.

\begin{table}[h]
\centering
\caption{Enceladus total tidal power against the closed-form synchronous
result, Eq.~\eqref{eq:closed}.}
\label{tab:enceladusclosed}
\begin{tabular}{lr}
\toprule
Source & Power (GW) \\
\midrule
\eqtide & 7.442710 \\
Eq.~\eqref{eq:closed}, computed independently & 7.442709 \\
Relative difference & $1.35\times10^{-7}$ \\
\bottomrule
\end{tabular}
\end{table}

\paragraph{Why no ocean column for Enceladus.} Equation~\eqref{eq:gamma} is a
\emph{surface}-ocean construction: it measures water depth against the solid
surface directly beneath it. Enceladus' ocean lies between a silicate interior
and a floating ice shell, so the geometry that makes $1+k_2-h_2$ the right
reduction does not apply, and the shell's own deformation --- not the sea
floor's --- sets the boundary the ocean responds to. Enceladus' $h_2$ has
never been measured in any case. We therefore report $k_2$, $h_2$ and
$\gamma_2$ for completeness but quote heights in Table~\ref{tab:enceladus} as
raising-potential values only. Supplying the correct reduction for a
subsurface ocean is the single largest piece of unfinished physics here, and
is the first item in \S\ref{sec:caveats}.

\section{Caveats and points for verification}
\label{sec:caveats}

The following are the places where independent scrutiny would be most useful.

\begin{enumerate}
\item \textbf{The subsurface-ocean reduction.} As above, $\gamma_2 = 1+k_2-h_2$
is derived for an ocean resting on a deforming solid surface. For an ocean
confined beneath a floating elastic shell the appropriate combination is
different, and depends on shell thickness and rigidity. Getting this right is
the prerequisite for using these amplitudes at Enceladus at all. \emph{This is
the item where we would most value a correction.}

\item \textbf{Amplitude normalization.} Equation~\eqref{eq:amp} uses a Kaula
normalization with every inclination function divided by 3. The icy-satellite
literature's familiar ``$3e$ radial, $4e$ librational'' shorthand is a
\emph{different} convention --- it is the coefficient of specific Legendre
terms --- so those numbers should not be compared directly without conversion.
Ratios within a species are exact under our convention. The frequencies and
powers carry no such ambiguity.

\item \textbf{Equilibrium tides only.} Even for Earth, Eq.~\eqref{eq:gamma}
gives the \emph{equilibrium} ocean tide. Real ocean tides depart from it by
large factors through basin resonance, coastline geometry, and self-attraction
and loading of the water itself, the last of which brings in the load Love
numbers $k_2'$ and $h_2'$ that are not modelled here.

\item \textbf{The default $h_2$.} Where $h_2$ is unmeasured we fall back on
$h_2 = \tfrac53 k_2$, exact for a homogeneous incompressible elastic sphere.
Earth shows the size of the error this can carry: $0.50$ against an observed
$0.61$. Any body with a differentiated interior or a decoupling ocean will
depart from it, plausibly by more.

\item \textbf{Dependence on the equilibrium-rotation convention.} The clean
collapse at Enceladus depends on taking the equilibrium spin to be exactly
synchronous, which holds in the discrete-lag treatment for $e < \sqrt{1/19}$.
Under the alternative convention the equilibrium rate is $(1+9.5e^2)n$ and the
semi-diurnal constituent moves off zero to a long-period line rather than a
permanent bulge. Which is physical depends on how one reads the CPL
truncation, and reasonable people may disagree.

\item \textbf{Enceladus $k_2$ and $Q$ are not well determined.} Tidal power
scales as $k_2/Q$; the values used ($0.1/20$) were chosen to land in the range
of the observed south-polar heat flow, roughly 5--16~GW. They are
illustrative, not a measurement. Reproducing the observed output from
equilibrium tidal heating alone is the well-known Enceladus heat-budget
problem and is \emph{not} claimed here.

\item \textbf{Truncation order.} The expansion is $\mathcal{O}(e)$ in
amplitude, hence $\mathcal{O}(e^2)$ in power, and first order in the
inclination functions. This is the order CPL itself retains; the constituent
decomposition inherits it and does not extend it. For Enceladus' $e=0.0047$
this is comfortable; for a high-eccentricity application it would not be.

\item \textbf{Obliquity in the Enceladus run} is set to a small nonzero value
($0.0005^\circ$) purely so the obliquity constituents are visible rather than
identically zero. Enceladus' obliquity is poorly constrained.
\end{enumerate}

\section*{Reproducing these results}

Every number in this note is produced by two example configurations
distributed with the code, one for each validation case, and each accompanied
by a script that regenerates its tables in the form shown here. Both cases are
also asserted as automated regression tests, covering the Doodson periods, the
frequency collapse, the zero-dissipation permanent tide, the exact rational
heating shares, the closed-form total, and the requirement that $h_2$ leave
the power budget untouched.

\end{document}
